\documentclass[aspectratio=169,11pt]{beamer}

% =========================
% Vietnamese Beamer Lecture
% =========================
% pdflatex-compatible Vietnamese setup
\usepackage[utf8]{inputenc}
\usepackage[T5]{fontenc}
\usepackage[vietnamese]{babel}
\usepackage{lmodern}
\usepackage{cmap}
\usepackage{amsmath,amssymb}
\usepackage{booktabs}
\usepackage{array}
\usepackage{tabularx}
\usepackage{listings}
\usepackage{xcolor}

\usetheme{Madrid}
\usecolortheme{seahorse}
\setbeamertemplate{navigation symbols}{}
\setbeamertemplate{footline}[frame number]

\definecolor{codebg}{RGB}{248,248,248}
\definecolor{keywordblue}{RGB}{0,70,140}
\definecolor{commentgray}{RGB}{90,90,90}

\lstdefinestyle{sqlstyle}{
  language=SQL,
  basicstyle=\ttfamily\scriptsize,
  keywordstyle=\color{keywordblue}\bfseries,
  commentstyle=\color{commentgray},
  backgroundcolor=\color{codebg},
  frame=single,
  framerule=0.2pt,
  breaklines=true,
  showstringspaces=false,
  columns=fullflexible
}

\title[Data Independence trong DBMS]{Bài giảng: Data Independence trong DBMS}
\subtitle{Độc lập dữ liệu trong hệ quản trị cơ sở dữ liệu}
\author{Biên soạn cho bài giảng DBMS}
\date{\today}

\AtBeginSection[]{
  \begin{frame}{Nội dung phần học}
    \tableofcontents[currentsection]
  \end{frame}
}

\begin{document}

%------------------------------------------------
\begin{frame}
  \titlepage
\end{frame}

%------------------------------------------------
\begin{frame}{Mục lục}
\begin{columns}[T,totalwidth=\textwidth]
  \column{0.48\textwidth}
  \tableofcontents[sections={1-4},hideallsubsections]
  \column{0.48\textwidth}
  \tableofcontents[sections={5-8},hideallsubsections]
\end{columns}
\end{frame}

%------------------------------------------------
\begin{frame}{Mục tiêu bài giảng}
Sau khi hoàn thành bài học này, người học có thể:
\begin{enumerate}
  \item Giải thích khái niệm \textbf{data independence} trong DBMS.
  \item Mô tả kiến trúc ba mức của DBMS: \textbf{Internal Level}, \textbf{Conceptual Level}, \textbf{View Level}.
  \item Phân biệt \textbf{Physical Data Independence} và \textbf{Logical Data Independence}.
  \item Phân tích ví dụ thực tế tương ứng với từng loại độc lập dữ liệu.
  \item Giải thích vai trò của data independence trong bảo trì và phát triển hệ thống.
\end{enumerate}
\end{frame}

%================================================
\section{Khái niệm Data Independence}

%------------------------------------------------
\begin{frame}{Data Independence là gì?}
\begin{block}{Định nghĩa}
\textbf{Data Independence} hay \textbf{độc lập dữ liệu} là khả năng thay đổi schema ở một mức trong kiến trúc cơ sở dữ liệu mà không ảnh hưởng đến schema ở mức cao hơn kế tiếp.
\end{block}

\vspace{0.3cm}
Nói cách khác:
\begin{itemize}
  \item Cách lưu trữ hoặc tổ chức dữ liệu bên trong có thể thay đổi.
  \item Người dùng và chương trình ứng dụng vẫn truy cập dữ liệu như trước.
  \item Mã ứng dụng không cần thay đổi đáng kể sau các tối ưu nội bộ.
\end{itemize}
\end{frame}

%------------------------------------------------
\begin{frame}{Ví dụ mở đầu}
Giả sử một hệ thống bán hàng thường xuyên tìm đơn hàng theo mã khách hàng.

\vspace{0.3cm}
DBA tạo thêm chỉ mục để tăng tốc truy vấn:
\begin{itemize}
  \item Cách lưu trữ hoặc truy cập dữ liệu bên trong thay đổi.
  \item Người dùng vẫn dùng cùng giao diện.
  \item Ứng dụng vẫn dùng cùng câu lệnh SQL.
\end{itemize}

\begin{alertblock}{Ý nghĩa}
Thay đổi bên trong database được DBMS che giấu khỏi ứng dụng và người dùng.
\end{alertblock}
\end{frame}

%------------------------------------------------
\begin{frame}{Vì sao Data Independence quan trọng?}
\begin{columns}[T]
\column{0.33\textwidth}
\begin{block}{Giảm bảo trì}
Developer không cần cập nhật ứng dụng sau mọi thay đổi nhỏ trong database.
\end{block}

\column{0.33\textwidth}
\begin{block}{Tăng linh hoạt}
Database có thể được tổ chức lại hoặc tối ưu nội bộ mà ít ảnh hưởng đến người dùng.
\end{block}

\column{0.33\textwidth}
\begin{block}{Phát triển lâu dài}
Database có thể mở rộng, cập nhật khi nghiệp vụ thay đổi mà không phá vỡ hệ thống hiện có.
\end{block}
\end{columns}
\end{frame}

%------------------------------------------------
\begin{frame}{Quiz nhanh: Khái niệm Data Independence}
\textbf{Câu 1.} Data independence trong DBMS là gì?
\begin{enumerate}[A.]
  \item Khả năng thay đổi schema ở một mức mà không ảnh hưởng đến mức cao hơn kế tiếp
  \item Khả năng xóa toàn bộ dữ liệu khỏi hệ thống
  \item Khả năng thay đổi màu giao diện người dùng
  \item Khả năng thay thế SQL bằng HTML
\end{enumerate}
\end{frame}

%------------------------------------------------
\begin{frame}{Quiz nhanh: Vai trò của Data Independence}
\textbf{Câu 2.} Data independence giúp giảm bảo trì vì lý do nào?
\begin{enumerate}[A.]
  \item Developer không cần sửa ứng dụng sau mọi thay đổi nhỏ trong database
  \item Người dùng phải tự quản lý ổ đĩa
  \item Database không bao giờ được thay đổi
  \item DBMS không còn cần schema
\end{enumerate}

\vspace{0.2cm}
\textbf{Câu 3.} Data independence hỗ trợ phát triển lâu dài như thế nào?
\begin{enumerate}[A.]
  \item Cho phép database được cập nhật và mở rộng mà ít làm hỏng hệ thống hiện có
  \item Bắt buộc viết lại toàn bộ ứng dụng mỗi năm
  \item Không cho phép thay đổi cấu trúc dữ liệu
  \item Chỉ dùng được với file văn bản
\end{enumerate}
\end{frame}

%================================================
\section{Kiến trúc ba mức của DBMS}

%------------------------------------------------
\begin{frame}{Kiến trúc ba mức của DBMS}
Để hiểu data independence, cần hiểu ba mức trong kiến trúc DBMS:

\vspace{0.4cm}
\begin{center}
\Large
\textbf{View Level}\\[0.2cm]
$\uparrow$\\[0.2cm]
\textbf{Conceptual Level}\\[0.2cm]
$\uparrow$\\[0.2cm]
\textbf{Internal Level}
\end{center}

\vspace{0.2cm}
\begin{alertblock}{Ý tưởng chính}
Mỗi mức che giấu chi tiết của mức thấp hơn để giảm sự phụ thuộc giữa ứng dụng và cách dữ liệu được lưu trữ.
\end{alertblock}
\end{frame}

%------------------------------------------------
\begin{frame}{Internal Level}
\begin{block}{Khái niệm}
\textbf{Internal Level} liên quan đến cách dữ liệu được lưu trữ vật lý trong hệ thống.
\end{block}

Mức này bao gồm:
\begin{itemize}
  \item File lưu trữ.
  \item Chỉ mục.
  \item Cách phân bổ không gian đĩa.
  \item Nén dữ liệu.
  \item Phương pháp truy cập dữ liệu.
  \item Cấu trúc lưu trữ bên trong.
\end{itemize}
\end{frame}

%------------------------------------------------
\begin{frame}{Conceptual Level}
\begin{block}{Khái niệm}
\textbf{Conceptual Level} mô tả cấu trúc logic của cơ sở dữ liệu.
\end{block}

Mức này bao gồm:
\begin{itemize}
  \item Bảng và trường dữ liệu.
  \item Thuộc tính.
  \item Quan hệ giữa các bảng.
  \item Khóa chính và khóa ngoại.
  \item Ràng buộc toàn vẹn.
\end{itemize}
\end{frame}

%------------------------------------------------
\begin{frame}{View Level}
\begin{block}{Khái niệm}
\textbf{View Level} mô tả cách người dùng và ứng dụng nhìn thấy dữ liệu.
\end{block}

Ví dụ trong hệ thống quản lý đào tạo:
\begin{itemize}
  \item Sinh viên chỉ xem điểm của chính mình.
  \item Giảng viên xem danh sách sinh viên trong lớp phụ trách.
  \item Phòng đào tạo xem toàn bộ dữ liệu học vụ.
  \item Phòng tài chính xem dữ liệu học phí.
\end{itemize}
\end{frame}

%------------------------------------------------
\begin{frame}{Mối liên hệ giữa ba mức}
\begin{tabularx}{\textwidth}{>{\bfseries}p{0.25\textwidth}X}
\toprule
Mức & Vai trò chính \\
\midrule
View Level & Cách từng nhóm người dùng hoặc ứng dụng nhìn thấy dữ liệu. \\
Conceptual Level & Mô hình logic tổng thể của database: bảng, cột, khóa, quan hệ. \\
Internal Level & Cách dữ liệu được lưu trữ và truy cập vật lý bên trong hệ thống. \\
\bottomrule
\end{tabularx}

\vspace{0.3cm}
\begin{alertblock}{Ghi nhớ}
Data independence xuất hiện khi thay đổi ở mức thấp không làm phá vỡ mức cao hơn.
\end{alertblock}
\end{frame}

%================================================
\section{Hai loại Data Independence}

%------------------------------------------------
\begin{frame}{Hai loại Data Independence}
\begin{tabularx}{\textwidth}{>{\bfseries}p{0.33\textwidth}p{0.25\textwidth}X}
\toprule
Loại độc lập dữ liệu & Thay đổi ở mức nào? & Không nên ảnh hưởng đến mức nào? \\
\midrule
Physical Data Independence & Internal Level & Conceptual Level và application programs \\
Logical Data Independence & Conceptual Level & View Level và application programs \\
\bottomrule
\end{tabularx}

\vspace{0.4cm}
\begin{block}{Cách hiểu nhanh}
\begin{itemize}
  \item \textbf{Physical}: thay đổi cách lưu trữ vật lý.
  \item \textbf{Logical}: thay đổi cấu trúc logic của dữ liệu.
\end{itemize}
\end{block}
\end{frame}

%================================================
\section{Physical Data Independence}

%------------------------------------------------
\begin{frame}{Physical Data Independence: Khái niệm}
\begin{block}{Định nghĩa}
\textbf{Physical Data Independence} là khả năng thay đổi cách dữ liệu được lưu trữ vật lý mà không ảnh hưởng đến schema logic hoặc ứng dụng người dùng.
\end{block}

\vspace{0.3cm}
Nó liên quan đến các thay đổi ở \textbf{Internal Level}.

\vspace{0.3cm}
\begin{alertblock}{Trọng tâm}
Tối ưu phần bên dưới database mà không bắt ứng dụng phía trên phải thay đổi.
\end{alertblock}
\end{frame}

%------------------------------------------------
\begin{frame}{Ví dụ thay đổi vật lý}
Các thay đổi sau thuộc Physical Data Independence:
\begin{itemize}
  \item Chuyển dữ liệu từ ổ C: sang ổ D:.
  \item Chuyển dữ liệu từ HDD sang SSD.
  \item Tạo index để tăng tốc truy vấn.
  \item Nén file dữ liệu để tiết kiệm dung lượng.
  \item Thay đổi access path hoặc phương pháp truy cập dữ liệu.
  \item Thay đổi cách tổ chức file lưu trữ.
\end{itemize}
\end{frame}

%------------------------------------------------
\begin{frame}[fragile]{Ví dụ SQL: Tạo index}
Giả sử hệ thống thường xuyên tìm đơn hàng theo \texttt{customer\_id}. DBA tạo index:

\begin{lstlisting}[style=sqlstyle]
CREATE INDEX idx_orders_customer
ON Orders(customer_id);
\end{lstlisting}

Ứng dụng vẫn có thể dùng truy vấn cũ:

\begin{lstlisting}[style=sqlstyle]
SELECT *
FROM Orders
WHERE customer_id = 101;
\end{lstlisting}

\begin{alertblock}{Nhận xét}
Schema logic của bảng \texttt{Orders} không thay đổi.
\end{alertblock}
\end{frame}

%------------------------------------------------
\begin{frame}{Lợi ích của Physical Data Independence}
\begin{itemize}
  \item Tối ưu backend mà không ảnh hưởng người dùng.
  \item Cải thiện hiệu năng truy vấn.
  \item Hỗ trợ nâng cấp phần cứng và lưu trữ.
  \item Giảm nhu cầu thay đổi application programs.
  \item Tăng khả năng bảo trì lâu dài.
\end{itemize}

\vspace{0.3cm}
\begin{block}{Thông điệp chính}
DBA có thể tối ưu hệ thống mà không làm gián đoạn ứng dụng đang chạy.
\end{block}
\end{frame}

%------------------------------------------------
\begin{frame}{Quiz: Physical Data Independence}
\textbf{Câu 1.} Physical Data Independence là gì?
\begin{enumerate}[A.]
  \item Khả năng thay đổi cách lưu trữ vật lý mà không ảnh hưởng đến schema logic hoặc ứng dụng
  \item Khả năng thay đổi giao diện người dùng
  \item Khả năng xóa toàn bộ database
  \item Khả năng thay đổi tên môn học
\end{enumerate}
\end{frame}

%------------------------------------------------
\begin{frame}{Quiz: Ví dụ và mục tiêu}
\textbf{Câu 2.} Ví dụ nào là Physical Data Independence?
\begin{enumerate}[A.]
  \item Tạo index để tăng tốc truy vấn mà ứng dụng không cần thay đổi
  \item Thêm trường email vào form nhập liệu
  \item Đổi màu website
  \item Xóa một chức năng khỏi giao diện
\end{enumerate}

\vspace{0.2cm}
\textbf{Câu 3.} Physical Data Independence thường phục vụ mục tiêu nào?
\begin{enumerate}[A.]
  \item Tối ưu hiệu năng và lưu trữ
  \item Thay đổi vai trò người dùng
  \item Thiết kế logo hệ thống
  \item Viết lại toàn bộ ứng dụng
\end{enumerate}
\end{frame}

%================================================
\section{Logical Data Independence}

%------------------------------------------------
\begin{frame}{Logical Data Independence: Khái niệm}
\begin{block}{Định nghĩa}
\textbf{Logical Data Independence} là khả năng thay đổi schema logic, tức là cấu trúc và quan hệ của dữ liệu, mà không ảnh hưởng đến view hoặc chương trình ứng dụng.
\end{block}

\vspace{0.3cm}
Nó liên quan đến các thay đổi ở \textbf{Conceptual Level}.

\vspace{0.3cm}
\begin{alertblock}{Trọng tâm}
Schema có thể phát triển theo nghiệp vụ mà vẫn giữ ổn định cho ứng dụng cũ.
\end{alertblock}
\end{frame}

%------------------------------------------------
\begin{frame}{Ví dụ thay đổi logic}
Các thay đổi sau thuộc Logical Data Independence:
\begin{itemize}
  \item Thêm cột mới vào bảng.
  \item Xóa cột không còn sử dụng.
  \item Tạo quan hệ mới giữa hai bảng.
  \item Tách một bảng thành nhiều bảng.
  \item Gộp nhiều bảng.
  \item Tạo view để đơn giản hóa truy cập.
\end{itemize}
\end{frame}

%------------------------------------------------
\begin{frame}[fragile]{Ví dụ SQL: Thêm cột mới}
Giả sử bảng \texttt{Employees} ban đầu có các cột:

\begin{lstlisting}[style=sqlstyle]
employee_id, full_name, department
\end{lstlisting}

Doanh nghiệp muốn lưu thêm email:

\begin{lstlisting}[style=sqlstyle]
ALTER TABLE Employees
ADD email VARCHAR(100);
\end{lstlisting}

Nếu các ứng dụng cũ chỉ dùng \texttt{employee\_id}, \texttt{full\_name}, \texttt{department}, chúng vẫn có thể hoạt động bình thường.
\end{frame}

%------------------------------------------------
\begin{frame}[fragile]{Vai trò của View}
Khi schema logic thay đổi, view có thể giúp giữ giao diện dữ liệu ổn định cho ứng dụng:

\begin{lstlisting}[style=sqlstyle]
CREATE VIEW EmployeeBasicView AS
SELECT employee_id, full_name, department
FROM Employees;
\end{lstlisting}

\begin{block}{Ý nghĩa}
Ứng dụng cũ truy vấn \texttt{EmployeeBasicView} mà không cần quan tâm bảng \texttt{Employees} đã có thêm cột \texttt{email}.
\end{block}
\end{frame}

%------------------------------------------------
\begin{frame}{Lợi ích của Logical Data Independence}
\begin{itemize}
  \item Dễ bảo trì mã ứng dụng hơn.
  \item Hỗ trợ cập nhật hệ thống khi nghiệp vụ phát triển.
  \item Giảm nhu cầu viết lại truy vấn cũ.
  \item Cho phép schema phát triển linh hoạt.
  \item Hỗ trợ nhiều view cho nhiều nhóm người dùng.
\end{itemize}

\vspace{0.3cm}
\begin{alertblock}{Lưu ý}
Logical Data Independence thường khó đạt được hơn Physical Data Independence vì schema logic gần với dữ liệu mà ứng dụng sử dụng.
\end{alertblock}
\end{frame}

%------------------------------------------------
\begin{frame}{Quiz: Logical Data Independence}
\textbf{Câu 1.} Logical Data Independence là gì?
\begin{enumerate}[A.]
  \item Khả năng thay đổi schema logic mà không ảnh hưởng đến view hoặc chương trình ứng dụng
  \item Khả năng thay đổi vị trí file vật lý
  \item Khả năng thay đổi ổ đĩa từ HDD sang SSD
  \item Khả năng thay đổi hệ điều hành
\end{enumerate}
\end{frame}

%------------------------------------------------
\begin{frame}{Quiz: Ví dụ và View}
\textbf{Câu 2.} Ví dụ nào là Logical Data Independence?
\begin{enumerate}[A.]
  \item Thêm cột \texttt{email} vào bảng \texttt{Employees} mà ứng dụng cũ vẫn hoạt động
  \item Tạo index trên ổ đĩa
  \item Chuyển dữ liệu sang SSD
  \item Nén file vật lý
\end{enumerate}

\vspace{0.2cm}
\textbf{Câu 3.} View hỗ trợ Logical Data Independence bằng cách nào?
\begin{enumerate}[A.]
  \item Giữ giao diện dữ liệu ổn định cho ứng dụng khi schema logic thay đổi
  \item Xóa toàn bộ dữ liệu cũ
  \item Thay đổi phần cứng máy chủ
  \item Không cho người dùng truy vấn database
\end{enumerate}
\end{frame}

%================================================
\section{So sánh và ví dụ tổng hợp}

%------------------------------------------------
\begin{frame}{So sánh Physical và Logical Data Independence}
\scriptsize
\begin{tabularx}{\textwidth}{>{\bfseries}p{0.22\textwidth}XX}
\toprule
Tiêu chí & Physical Data Independence & Logical Data Independence \\
\midrule
Trọng tâm & Cách dữ liệu được lưu trữ vật lý & Cấu trúc và tổ chức dữ liệu \\
Mức liên quan & Internal Schema & Conceptual Schema \\
Loại thay đổi & File, index, ổ đĩa, nén, access path & Bảng, cột, quan hệ, view \\
Ảnh hưởng đến ứng dụng & Không nên ảnh hưởng đến application programs & Có thể khó hơn để tránh ảnh hưởng \\
Mục tiêu & Tối ưu hiệu năng và lưu trữ & Phát triển thiết kế database \\
Mức độ đạt được & Dễ hơn & Khó hơn \\
Ví dụ & Di chuyển file dữ liệu hoặc thêm index & Thêm/xóa cột trong bảng \\
\bottomrule
\end{tabularx}
\end{frame}

%------------------------------------------------
\begin{frame}[fragile]{Ví dụ tổng hợp: Hệ thống quản lý bán hàng}
\textbf{Physical Data Independence:} DBA tạo index trên cột \texttt{customer\_id} trong bảng \texttt{Orders}.

\begin{lstlisting}[style=sqlstyle]
CREATE INDEX idx_customer_id
ON Orders(customer_id);

SELECT *
FROM Orders
WHERE customer_id = 1001;
\end{lstlisting}

\begin{block}{Nhận xét}
Cách truy cập dữ liệu được tối ưu, nhưng schema logic và ứng dụng không đổi.
\end{block}
\end{frame}

%------------------------------------------------
\begin{frame}[fragile]{Ví dụ tổng hợp: Thay đổi schema logic}
\textbf{Logical Data Independence:} Doanh nghiệp muốn lưu thêm email của khách hàng.

\begin{lstlisting}[style=sqlstyle]
ALTER TABLE Customers
ADD email VARCHAR(100);
\end{lstlisting}

Các báo cáo cũ chỉ dùng \texttt{customer\_id} và \texttt{customer\_name} vẫn hoạt động.

\begin{lstlisting}[style=sqlstyle]
CREATE VIEW CustomerBasicView AS
SELECT customer_id, customer_name
FROM Customers;
\end{lstlisting}
\end{frame}

%================================================
\section{Câu hỏi ôn tập và bài tập}

%------------------------------------------------
\begin{frame}{Câu hỏi trắc nghiệm ôn tập 1}
\small
\textbf{Câu 1.} Data independence là gì?
\begin{enumerate}[A.]
  \item Khả năng thay đổi schema ở một mức mà không ảnh hưởng đến mức cao hơn kế tiếp
  \item Khả năng xóa dữ liệu không kiểm soát
  \item Khả năng đổi màu giao diện
  \item Khả năng thay thế database bằng file ảnh
\end{enumerate}

\textbf{Câu 2.} Có mấy loại data independence chính?
\begin{enumerate}[A.]
  \item 1
  \item 2
  \item 3
  \item 4
\end{enumerate}
\end{frame}

%------------------------------------------------
\begin{frame}{Câu hỏi trắc nghiệm ôn tập 2}
\small
\textbf{Câu 3.} Physical Data Independence liên quan chủ yếu đến mức nào?
\begin{enumerate}[A.]
  \item Internal Level
  \item View Level
  \item User Interface
  \item Application Screen
\end{enumerate}

\textbf{Câu 4.} Logical Data Independence liên quan chủ yếu đến mức nào?
\begin{enumerate}[A.]
  \item Internal Level
  \item Conceptual Level
  \item Disk Level
  \item Hardware Level
\end{enumerate}
\end{frame}

%------------------------------------------------
\begin{frame}{Câu hỏi trắc nghiệm ôn tập 3}
\small
\textbf{Câu 5.} Ví dụ nào là Physical Data Independence?
\begin{enumerate}[A.]
  \item Tạo index mới mà không làm thay đổi ứng dụng
  \item Thêm trường email vào giao diện
  \item Đổi màu trang web
  \item Xóa một nút trên màn hình
\end{enumerate}

\textbf{Câu 6.} Ví dụ nào là Logical Data Independence?
\begin{enumerate}[A.]
  \item Thêm cột mới vào bảng mà ứng dụng cũ vẫn hoạt động
  \item Chuyển dữ liệu sang SSD
  \item Tạo file vật lý mới
  \item Thay đổi block size trên đĩa
\end{enumerate}
\end{frame}

%------------------------------------------------
\begin{frame}{Câu hỏi trắc nghiệm ôn tập 4}
\small
\textbf{Câu 7.} Physical Data Independence chủ yếu dùng để làm gì?
\begin{enumerate}[A.]
  \item Tối ưu hiệu năng và lưu trữ
  \item Thay đổi màu giao diện
  \item Xóa dữ liệu người dùng
  \item Thay thế DBMS
\end{enumerate}

\textbf{Câu 8.} Logical Data Independence thường khó hơn Physical Data Independence vì sao?
\begin{enumerate}[A.]
  \item Vì thay đổi logic thường gần với cấu trúc mà ứng dụng đang sử dụng
  \item Vì thay đổi vật lý luôn làm hỏng ứng dụng
  \item Vì view không liên quan đến dữ liệu
  \item Vì logical schema không tồn tại
\end{enumerate}
\end{frame}

%------------------------------------------------
\begin{frame}{Câu hỏi trắc nghiệm ôn tập 5}
\small
\textbf{Câu 9.} View có thể giúp gì cho Logical Data Independence?
\begin{enumerate}[A.]
  \item Giữ giao diện dữ liệu ổn định cho ứng dụng
  \item Xóa database
  \item Di chuyển file sang ổ cứng khác
  \item Thay đổi RAM máy chủ
\end{enumerate}

\textbf{Câu 10.} Data independence giúp ứng dụng như thế nào?
\begin{enumerate}[A.]
  \item Giảm khả năng bị ảnh hưởng bởi thay đổi trong database
  \item Bắt buộc viết lại ứng dụng sau mọi thay đổi
  \item Không cho phép thay đổi schema
  \item Xóa toàn bộ view
\end{enumerate}
\end{frame}

%------------------------------------------------
\begin{frame}{Câu hỏi tự luận ngắn}
\begin{enumerate}
  \item Giải thích khái niệm data independence trong DBMS.
  \item Phân biệt Physical Data Independence và Logical Data Independence.
  \item Vì sao Physical Data Independence giúp tối ưu hiệu năng mà không ảnh hưởng đến ứng dụng?
  \item Vì sao Logical Data Independence thường khó đạt được hơn Physical Data Independence?
  \item Hãy giải thích vai trò của view trong việc hỗ trợ Logical Data Independence.
\end{enumerate}
\end{frame}

%------------------------------------------------
\begin{frame}{Bài tập vận dụng 1}
Một DBA chuyển dữ liệu từ ổ HDD sang SSD và tạo thêm index để tăng tốc truy vấn. Ứng dụng không cần thay đổi.

\vspace{0.4cm}
\begin{block}{Yêu cầu}
Hãy xác định đây là loại data independence nào và giải thích.
\end{block}
\end{frame}

%------------------------------------------------
\begin{frame}{Bài tập vận dụng 2}
Một hệ thống quản lý nhân sự thêm cột \texttt{email} vào bảng \texttt{Employees}. Các báo cáo cũ vẫn chỉ hiển thị mã nhân viên, tên và phòng ban.

\vspace{0.4cm}
\begin{block}{Yêu cầu}
Hãy xác định đây là loại data independence nào và giải thích.
\end{block}
\end{frame}

%------------------------------------------------
\begin{frame}{Bài tập vận dụng 3}
Một hệ thống bán hàng muốn thay đổi cấu trúc bảng \texttt{Customers}, nhưng vẫn muốn ứng dụng cũ tiếp tục truy vấn \texttt{customer\_id} và \texttt{customer\_name}.

\vspace{0.4cm}
\begin{block}{Yêu cầu}
Hãy đề xuất cách dùng view để giảm ảnh hưởng đến ứng dụng.
\end{block}
\end{frame}

%------------------------------------------------
\begin{frame}{Bài tập vận dụng 4}
\begin{block}{Yêu cầu}
Hãy cho một ví dụ trong hệ thống quản lý sinh viên minh họa cả:
\begin{itemize}
  \item Physical Data Independence.
  \item Logical Data Independence.
\end{itemize}
\end{block}

\vspace{0.3cm}
Gợi ý: Có thể xét các thay đổi về index, lưu trữ, bảng \texttt{Students}, hoặc view cho từng nhóm người dùng.
\end{frame}

%================================================
\section{Tóm tắt}

%------------------------------------------------
\begin{frame}{Tóm tắt bài học}
\begin{itemize}
  \item Data independence là khả năng thay đổi schema ở một mức mà không ảnh hưởng đến mức cao hơn kế tiếp.
  \item DBMS thường có ba mức: Internal Level, Conceptual Level và View Level.
  \item Physical Data Independence cho phép thay đổi cách lưu trữ vật lý mà không ảnh hưởng đến schema logic hoặc ứng dụng.
  \item Logical Data Independence cho phép thay đổi schema logic mà không ảnh hưởng đến view hoặc chương trình ứng dụng.
  \item Physical Data Independence thường dễ đạt được hơn Logical Data Independence.
  \item View có thể giúp che giấu thay đổi ở schema logic và giữ ứng dụng ổn định.
\end{itemize}
\end{frame}

%------------------------------------------------
\begin{frame}{Từ khóa chính}
\begin{columns}[T]
\column{0.5\textwidth}
\begin{itemize}
  \item Data Independence
  \item Physical Data Independence
  \item Logical Data Independence
  \item DBMS
  \item Internal Level
  \item Conceptual Level
  \item View Level
\end{itemize}

\column{0.5\textwidth}
\begin{itemize}
  \item Internal Schema
  \item Conceptual Schema
  \item View
  \item Index
  \item Schema
  \item Database Administrator
  \item Application Program
\end{itemize}
\end{columns}
\end{frame}

%------------------------------------------------
\begin{frame}{Kết thúc bài học}
\begin{center}
\Huge Câu hỏi và thảo luận
\end{center}
\end{frame}

\end{document}
